% !TEX root = project_report.tex
\documentclass[11pt,a4paper]{article}
\usepackage[margin=22mm,headheight=15pt]{geometry}
\usepackage{fontspec}
\IfFontExistsTF{Liberation Serif}{\setmainfont{Liberation Serif}}{\IfFontExistsTF{Times New Roman}{\setmainfont{Times New Roman}}{\setmainfont{Latin Modern Roman}}}
\IfFontExistsTF{Liberation Sans}{\setsansfont{Liberation Sans}}{\IfFontExistsTF{Arial}{\setsansfont{Arial}}{\setsansfont{Latin Modern Sans}}}
\IfFontExistsTF{Liberation Mono}{\setmonofont{Liberation Mono}[Scale=MatchLowercase]}{\IfFontExistsTF{Courier New}{\setmonofont{Courier New}[Scale=MatchLowercase]}{\setmonofont{Latin Modern Mono}[Scale=MatchLowercase]}}
\usepackage{graphicx,booktabs,array,amsmath,xcolor,hyperref,fancyhdr,caption,titlesec,microtype,tabularx,placeins}
\definecolor{ink}{HTML}{213B48}
\definecolor{teal}{HTML}{086D78}
\definecolor{muted}{HTML}{596A73}
\hypersetup{colorlinks=true,linkcolor=teal,urlcolor=teal,citecolor=teal,
  pdftitle={Rainfall across Singapore: Patterns and next-30-minute risk},
  pdfsubject={DSA5208 group project, 2017--2024 rainfall observations}}
\urlstyle{same}
\setlength{\parindent}{0pt}
\setlength{\parskip}{5pt plus 1pt minus 1pt}
\setlength{\textfloatsep}{13pt plus 2pt minus 2pt}
\setlength{\intextsep}{13pt plus 2pt minus 2pt}
\setlength{\emergencystretch}{2em}
\linespread{1.025}
\raggedbottom
\renewcommand{\topfraction}{.90}
\renewcommand{\bottomfraction}{.80}
\renewcommand{\textfraction}{.08}
\renewcommand{\floatpagefraction}{.75}
\setcounter{topnumber}{3}
\setcounter{bottomnumber}{2}
\setcounter{totalnumber}{5}
\titleformat{\section}{\large\sffamily\bfseries\color{ink}}{\thesection}{.65em}{}
\titleformat{\subsection}{\normalsize\sffamily\bfseries\color{ink}}{\thesubsection}{.65em}{}
\titlespacing*{\section}{0pt}{14pt}{6pt}
\titlespacing*{\subsection}{0pt}{10pt}{4pt}
\captionsetup{font=small,labelfont=bf,labelsep=period,justification=raggedright,singlelinecheck=false,skip=5pt}
\pagestyle{fancy}\fancyhf{}
\fancyhead[L]{\footnotesize\sffamily\color{muted}Rainfall across Singapore}
\fancyhead[R]{\footnotesize\sffamily\color{muted}2017--2024}
\fancyfoot[C]{\footnotesize\thepage}
\renewcommand{\headrulewidth}{.25pt}
\graphicspath{{../figures/project-report/}{../figures/task1-v5-report/}{../figures/task2-report/}}
\newcommand{\plot}[3]{\par\begin{figure}[htbp]\centering\includegraphics[width=\linewidth]{#1.pdf}\caption{#2}\label{#3}\end{figure}\par}
\begin{document}
\thispagestyle{empty}
{\small\sffamily\color{teal}DSA5208 GROUP PROJECT}\par
\vspace{3mm}
{\LARGE\sffamily\bfseries\color{ink}Rainfall across Singapore}\par
\vspace{1mm}
{\large\sffamily Patterns and next-30-minute risk, 2017--2024}\par
\vspace{4mm}
\begin{tabularx}{\linewidth}{@{}>{\raggedright\arraybackslash}X>{\raggedright\arraybackslash}X@{}}
\small Group member name & \small Student ID\\[5pt]
Ma Xiaoxiao & A0259173X\\[5pt]
Ye Sifan & A0262394Y\\[5pt]
Chen Yuda & A0353182H\\
\end{tabularx}
\vspace{2mm}
% !TEX root = ../project_report.tex
\begin{abstract}
Singapore's five-minute rainfall records offer a detailed view of local weather, while uneven coverage affects comparisons over time. We analyse the 2017--2024 annual datasets with Apache Spark, retaining 48.54 million observations from 91 station IDs after cleaning. Among sufficiently complete station-months, November has the highest mean monthly rainfall and wet intervals are most frequent in the afternoon. For prediction, we compare random forests, gradient-boosted trees, logistic regression and historical frequencies under two information settings: recent station rainfall and station/calendar information alone. On 5.73 million shared 2024 test cases, the recent-rainfall forest achieves mean Brier 0.01613, compared with 0.02283 for the calendar forest. RF has the lowest observed mean Brier among the tested classifier pipelines; historical frequency performs best for calendar-only risk. A paired day-block analysis supports a small RF advantage over GBT with recent rainfall, while RF fitting fractions from 15\% to 100\% show no clear mean improvement. Gains are concentrated at low rainfall thresholds, and heavy-rain probability skill remains limited. Later experiments reuse 2024, so model selection remains exploratory. These results support the recent-rainfall RF as a reference model and motivate stronger validation and tail calibration.
\end{abstract}

\section{Introduction}
The project asks two related questions: how rainfall varies across Singapore within and between years, and how well we can estimate a station's chance of receiving more than a specified amount in the next 30 minutes. The first question concerns patterns across many observations. The second requires a clear account of what information would be available when a forecast is issued.

We use all eight annual datasets from 2017 to 2024. For exploration, the main difficulty is uneven coverage: a low recorded monthly total can reflect missing observations as well as dry weather. We address this through completeness rules and comparisons at common stations. For prediction, we distinguish a calendar-only model from a model that also uses recently published rainfall. This lets us examine both the three-input query in the assignment---station, prediction time and threshold---and the benefit of supplying current weather observations internally.

The report follows these two questions from data preparation to results. Spark handles the observation-level calculations and model fitting; figures and supplementary analyses use the saved result tables.

% !TEX root = ../project_report.tex
\section{Data and methods}

\subsection{What the records contain}
We use the eight annual Historical Rainfall across Singapore datasets for 2017--2024. Each observation contains a station identifier, coordinates, a timestamp, rainfall in millimetres and update metadata. All times are interpreted in Singapore time. The measurement is an amount accumulated over a five-minute interval, rather than a rainfall rate. The publisher cautions that records may be missing and have not undergone the quality control applied to official climate records \cite{data}.

The cleaned dataset contains 48,538,625 observations from 91 station IDs. Recorded zeros account for 97.08\% of observations, while positive amounts have a long upper tail (Figure~\ref{fig:profile}). We therefore examine how often rain occurs separately from how much falls during wet intervals. A zero is a measurement of no rain; a missing record provides no such information. Appendix~\ref{app:data} provides the positive-rainfall distribution and a raw-versus-clean source-gap example.

\plot{01_data_profile}{The observation sample by year: record counts, the percentage of recorded zeros, and quantiles of positive five-minute rainfall. Positive-rainfall quantiles exclude dry intervals.}{fig:profile}

\subsection{Cleaning and missing observations}
The PySpark pipeline standardises identifiers and field types, checks rainfall units and measurement type, and rejects malformed records, invalid timestamps, invalid coordinates, and non-finite or negative rainfall. Timestamps within two seconds of a five-minute boundary are rounded to that boundary; records farther from the grid are excluded. Exact duplicates are removed. If records sharing a station and time disagree on rainfall or coordinates, or include an invalid source record, all records for that station and time are set aside for inspection. Consistent measurements with different update metadata are consolidated using the latest publication timestamps, conservatively delaying their availability for prediction.

The source-to-clean audit reconciles 48,538,713 raw rows with 48,538,625 clean observations. The count falls by 88; because cleaning also consolidates repeated records, this difference is not a count of invalid measurements. All raw timestamps were parseable. High non-negative amounts are retained for review, including the maximum of 37.2 mm in five minutes; the cleaning rules alone do not establish whether an extreme reading is meteorologically valid. The pipeline preserves the raw files and writes rejected records, quality summaries and clean Parquet partitioned by year and month. The repository includes the summary evidence; the full raw and intermediate datasets require reproduction through the documented cloud workflow.

Missingness matters more to this analysis than the small change in row count. The source audit identifies 63 calendar dates without observations anywhere in the network. August 2018, for example, contains 36,410 raw records, all within its first three days. This gap was already present in the annual source file. Timestamp rounding can also move near-midnight records to an adjacent date, so daily raw-to-clean differences need not represent deletions.

For station $s$, year $y$ and month $m$, we measure completeness as
\[
 C_{sym}=\frac{N_{sym}}{288D_{ym}},
\]
where $N_{sym}$ is the number of valid unique observations and $D_{ym}$ is the number of calendar days. The analysis constructs a 91-station by 96-month grid; absent months have zero observations and an undefined rainfall total. This grid does not imply that every station operated throughout the study. Low coverage can reflect both operating periods and missing records.

\plot{03_coverage}{Monthly data availability. The upper panel gives median coverage among stations with any observations; the lower panel counts stations meeting the 90\% rule. Grey cells in the lower panel have no qualifying stations.}{fig:coverage}

\subsection{Comparing rainfall over time}
Monthly summaries require at least 90\% coverage. Annual summaries require 95\% coverage overall and 90\% in every month, preventing a concentrated seasonal gap from being hidden by a high annual average. Daily-distribution summaries require 95\% daily coverage within eligible months. These rules limit missing-data exposure but cannot remove bias if missingness is related to weather. Totals remain observed amounts, without gap filling or scaling by inverse coverage.

Within each year-month, rainfall amount gives eligible station totals equal weight. Wet frequency and conditional intensity instead pool eligible observed intervals: wet frequency is the number of positive intervals divided by observed intervals, and conditional intensity is their rainfall sum divided by positive intervals. Seasonal curves then give available years equal weight. Hourly summaries pool eligible observations, giving more weight to periods with more observed intervals. Station means describe the sampled network, rather than an area-weighted national rainfall total.

Cross-year comparisons match station IDs within the same calendar month for each year pair. Coordinate stability screens use a tolerance of 0.005 degrees per axis within and between paired periods. We average the within-station changes and express percentage change relative to the earlier pair mean. Plots show only comparisons with at least five common stations; this is a display threshold, not a significance test. Different year pairs can contain different stations.

\subsection{Computing environment}
Observation-level processing, feature construction and Spark MLlib training ran on Google Cloud Dataproc. Table~\ref{tab:environment} records the experimental configuration. Only saved summaries were downloaded for local analysis and plotting; source snapshots and job metadata accompany the results.

\begin{table}[htbp]\small\centering
\begin{tabularx}{\linewidth}{>{\raggedright\arraybackslash}p{.22\linewidth}>{\raggedright\arraybackslash}X}
\toprule Component & Configuration\\\midrule
Cluster & One master and two workers in \texttt{asia-southeast1-c}\\
Each VM & \texttt{e2-standard-4}: 4 vCPUs, 16 GB RAM, 80 GB \texttt{pd-balanced} boot disk\\
Worker capacity & 8 vCPUs and 32 GB RAM, excluding the master\\
Runtime & Dataproc 2.2.87, Debian 12; Spark/PySpark 3.5.3 on YARN\\
Spark settings & 6 GB driver memory; 1 GB maximum driver result; 128 shuffle partitions; adaptive SQL; Singapore timezone\\
Local figures & Apple M2, 16 GB RAM; pandas, Matplotlib and Plotly\\
Storage & Cloud Storage for raw files, Parquet data and result tables\\
\bottomrule
\end{tabularx}
\caption{Configuration used for the reported experiments.}
\label{tab:environment}
\end{table}

\FloatBarrier
% !TEX root = ../project_report.tex
\section{Rainfall patterns, 2017--2024}

\subsection{Seasonality and changes between years}

The clearest seasonal feature is the increase in rainfall towards the end of the year. November has the highest mean observed monthly amount, 320.8~mm, while February has the lowest, 146.1~mm. These averages give equal weight to each available year after averaging its qualified stations. Figure~\ref{fig:monthly} shows the individual year-months behind this summary; Figure~\ref{fig:seasonal} separates rainfall amount into frequency and intensity. The broad pattern agrees with the wetter November--January period and relatively dry February described by the Meteorological Service Singapore \cite{climate}.

\plot{05_monthly_amount}{Mean observed monthly rainfall at qualified stations, 2017--2024. All cells share one colour scale. Grey cells have no qualified station-months.}{fig:monthly}

The frequency and intensity curves help explain what makes a month wet. November has the highest mean wet-interval frequency, approximately 4.30\%, but its mean amount per wet interval is about 0.89~mm. March has a higher conditional intensity, about 1.06~mm, despite a lower monthly total. Thus, November's high accumulation is associated with more frequent rain, rather than the largest amount in each wet interval. This distinction matters for prediction: a season with frequent rain need not have the greatest risk of a large amount within a short window.

\plot{06_seasonality}{Seasonal means of monthly rainfall, wet-interval frequency and rainfall per wet interval. Available years receive equal weight; shading shows their minimum--maximum range.}{fig:seasonal}

Monthly totals also reflect the number of days in a month. February's shorter duration contributes to its low total, so the amount curve alone cannot establish that it has the lowest propensity for rain. The frequency and intensity panels provide complementary measures. Between six and eight years contribute to each calendar-month average, and the shaded ranges describe their observed variation rather than uncertainty in a long-term climate normal.

\begingroup\clubpenalty=10000
To examine changes between years, we compare the same calendar month at stations present in both periods and meeting the location-stability checks. December 2024, for example, had a mean of 213.9~mm across 53 common stations, compared with 379.8~mm in December 2023: a reduction of 43.7\%. This provides stronger evidence of a change at those locations than a comparison of two unmatched network means. Figure~\ref{fig:paired} applies the same approach to adjacent-year monthly pairs. The participating stations vary between pairs, so the percentages describe separate comparisons and cannot be chained into a national trend.\par\endgroup

\plot{09_paired_changes}{Adjacent-year changes in monthly rainfall at common stable-location stations, relative to the earlier year's mean. Grey cells lack support from at least five common stations; this cutoff is a display rule, not a significance test.}{fig:paired}

Full-year comparisons are more limited. Only 2017, 2022 and 2024 contain stations meeting both the annual and every-month completeness requirements. Their mean observed annual totals are 2,396.6~mm (9 stations), 2,920.9~mm (54 stations) and 2,682.3~mm (29 stations), respectively. Among the 26 stations eligible in both 2022 and 2024, the mean declined by 284.8~mm, or 9.6\%. Comparisons involving 2017 have only three or four common stations. These results support a limited comparison of monitored locations, but neither an eight-year annual trend nor a national wettest-year ranking. Months with adequate coverage in the other years remain part of the monthly analysis.

\subsection{When and where rain occurs}

Rain is more frequent in the afternoon. Across valid intervals in qualified station-months, wet frequency reaches 5.80\% at 15:00--15:59 and falls to 1.10\% at 23:00--23:59, using Singapore time. Figure~\ref{fig:diurnal} shows how this daily pattern varies by month and year. The afternoon enhancement is consistent with MSS's description of Singapore's rainfall \cite{climate} and motivates including hour-of-day features in the predictive model. These pooled frequencies describe observed occurrence; they are not forecasts for a particular afternoon, and the yearly curves can contain different mixes of months and stations.

\plot{10_diurnal}{Wet-observation frequency by month and hour, and by year and hour, in Singapore time. Frequencies are weighted by the number of observed intervals.}{fig:diurnal}

Rainfall also varies substantially across monitoring locations. In December 2024, qualified station totals range from 29.5~mm at S224 to 326.8~mm at S117. The December maps retain 41--69 stations in each of the eight years, making this month a useful illustration of spatial and year-to-year variation. Figure~\ref{fig:spatialdec} shows the measured station values on a common scale. The full set of 96 year-month maps remains in the project repository's supplementary gallery. The map outline comes from the retained boundary assets \cite{geo}. We avoid interpolating a continuous surface because the station network alone does not establish how rainfall varies between gauges; station means likewise should not be interpreted as area-weighted national rainfall.

\plot{13_spatial_month_12}{December rainfall at qualified stations in each study year. All eight panels use a 0--600 mm scale, covering the full range of December station totals. Panel labels give station counts; colours represent measurements at gauges, with no spatial interpolation.}{fig:spatialdec}

\subsection{How sensitive are these summaries to coverage?}

Changing the monthly coverage cutoff produces relatively small changes in the supported monthly means. Relative to the 90\% cutoff, using 80\% gives a median absolute change of 0.61\% across 85 shared year-months, with a maximum of 5.78\%. Using 95\% gives a median change of 0.40\% across 84 shared year-months, with a maximum of 2.97\%. This is reassuring for the descriptive amount summaries, although each cutoff changes the contributing stations. It does not establish that missing observations are random. Residual gaps, uneven station coverage and the short study period remain the main reasons to treat these findings as patterns in the available observations rather than evidence of long-term climate change.

Appendices~\ref{app:temporal}--\ref{app:spatial} retain monthly trajectories, annual station totals, daily-rainfall summaries, coverage-sensitivity curves and additional spatial views. They supplement the principal comparisons without requiring all stations to qualify in all eight years.

\FloatBarrier
% !TEX root = ../project_report.tex
\section{Predicting the next 30 minutes}
\subsection{Prediction target and available information}
For station $s$, issue time $t$ and threshold $q$ in millimetres, the target is
\[
Y_{s,t}=\sum_{j=1}^{6}r_{s,t+5j\,\mathrm{min}},\qquad
p_{s,t}(q)=P(Y_{s,t}>q\mid X_{s,t}).
\]
The output is the probability of exceeding the requested amount. Observation timestamps mark the assumed endpoints of five-minute accumulation intervals. Labels require all six exact future observations; totals are rounded to 0.001 mm and exceedance is strictly greater than $q$. Examples with missing future observations are excluded.

\textbf{A (short-term forecasting)} uses recent station rainfall as well as station and calendar features. \textbf{B (calendar risk)} uses only station and calendar features to estimate historical risk at a requested time. We fit the same model families in both settings to compare short-term forecasting with calendar risk. The query's three user inputs are station, issue time and threshold; A derives additional features internally from the supplied observation archive.

A reads a supplied clean observation table and requires twelve history endpoints at $t-65,t-60,\ldots,t-10$ minutes. Both recorded update timestamps for every history record must be no later than $t$. The ten-minute buffer reflects the typical publication delay (Table~\ref{tab:updatedelay}); the record-level checks also exclude observations with longer delays. Missing or unavailable history rejects the example or query. Figure~\ref{fig:t2timeline} separates this history from the future target.

\begin{table}[htbp]\centering\small
\begin{tabular}{lrrrr}\toprule
Period & Observations & Median (min) & P90 (min) & Within 10 min (\%)\\\midrule
2017--2021 (training) & 28,256,245 & 9.58 & 15.08 & 68.93 \\
2017--2024 (all years) & 48,538,625 & 9.47 & 12.68 & 76.40 \\
\bottomrule\end{tabular}
\caption{Observation-to-update delay in the cleaned records, using the later of the two update timestamps. P90 is the 90th percentile; the final column gives the percentage updated within ten minutes.}
\label{tab:updatedelay}
\end{table}

\plot{17_prediction_timeline}{The information used by A at issue time $t$. Markers show interval endpoints. The latest history interval ends ten minutes before issue time; the six intervals after issue time form the target.}{fig:t2timeline}
\FloatBarrier
\subsection{Features and their purpose}
Location and station ID represent spatial differences; cyclic time encodings represent the seasonal and daily patterns explored in Part 1 without a discontinuity at midnight or New Year. Annual phase accounts for leap years, and numeric year is excluded. A's rainfall summaries represent recent intensity, persistence and change; all history windows end at $t-10$ minutes (Table~\ref{tab:input_features}).
\begin{table}[htbp]\centering\small
\begin{tabular}{p{.63\linewidth}cc}\toprule
Feature group (number of numeric features) & A & B\\\midrule
Latitude and longitude (2) & Yes & Yes\\
Time of day: sine and cosine (2) & Yes & Yes\\
Annual phase: sine and cosine (2) & Yes & Yes\\
Latest three five-minute rainfall amounts (3) & Yes & No\\
Rainfall totals over 15, 30 and 60 minutes (3) & Yes & No\\
Maximum five-minute rainfall within 30 minutes (1) & Yes & No\\
Wet-interval count and dry-interval fraction over 60 minutes (2) & Yes & No\\
Difference between the latest and preceding 15-minute totals (1) & Yes & No\\\midrule
\textbf{Total numeric features} & \textbf{16} & \textbf{6}\\
Station ID (categorical; counted separately) & Yes & Yes\\\bottomrule
\end{tabular}
\caption{Input features for the two information settings. A uses recent rainfall in addition to the station and calendar inputs shared with B. For LR and GBT, the query threshold $q$ is an additional input; RF uses it when converting class probabilities to exceedance probabilities.}
\label{tab:input_features}
\end{table}

\FloatBarrier
\subsection{Models: answering a query at any rainfall threshold}
We compare a historical-frequency baseline with three Spark MLlib classifiers~\cite{sparkml}. The key modelling choice is how to turn their outputs into $P(Y>q)$ for a user-supplied threshold, including values absent from the fitting grid.

\textbf{Historical frequency} estimates the exceedance rate for each station/month/hour group, shrinking it towards the global training rate using 100 pseudo-observations. Unseen groups use the global rate. This establishes how much can be achieved from recurring rainfall patterns before adding a classifier.

\textbf{Distributional random forest (RF)} predicts seven amount bins: $\{0\}$, $(0,1]$, $(1,2]$, $(2,5]$, $(5,10]$, $(10,20]$ and $(20,\infty)$. Predicting a distribution lets one fitted forest answer different rainfall thresholds. For a query inside a bin, we use the observed training distribution within that bin:

\[
\widehat p(q\mid X)=\sum_{k=0}^{6}\widetilde\pi_k(X)\widehat S_k(q),
\qquad
\widetilde\pi_k=\frac{\pi_k^{1/T}}{\sum_\ell\pi_\ell^{1/T}}.
\]
Here $\pi_k(X)$ is the forest probability for bin $k$, and $\widehat S_k(q)$ is the fraction of training amounts in that bin exceeding $q$. The temperature $T$ adjusts how concentrated the bin probabilities are; $T=1$ leaves them unchanged. The resulting exceedance probability is non-increasing in the requested rainfall threshold. The within-bin distributions are shared across weather conditions and extend to the largest training total, 71.8 mm; the model returns zero at or above this empirical maximum. This is a finite-support limitation, not a claim that larger future accumulations are impossible. The query output flags thresholds outside the training support.

Station IDs are encoded as categorical features, and numeric imputation uses training-set medians. The forests in the main model/fraction comparison use 40 trees, depth 8, a minimum of 25 examples per child node, square-root feature selection, subsampling 0.8 and at most 90 feature bins. The 15\% and 100\% searches select depth 8 from 5 and 8; the 30\% comparison reuses the same reference configuration. Temperature selection chooses $T=1$ for all three fractions in both settings on the 2023 calibration set.

\textbf{Threshold-conditioned logistic regression (LR)} instead treats $q$ as an input and predicts the binary label $\mathbf1(Y>q)$. Each sampled station-time example is expanded to seven thresholds: 0, 0.2, 1, 2, 5, 10 and 20 mm. Station IDs are one-hot encoded and numeric inputs scaled. L2 regularisation is selected from 0.01 and 0.1, with at most 100 optimisation iterations. LR supplies a simple direct-probability comparator and receives no separate calibration fit.

\textbf{Threshold-conditioned gradient-boosted trees (GBT).} GBT uses the same binary labels and threshold expansion as LR, but learns nonlinear relationships by fitting trees sequentially. Its pipeline indexes station IDs, imputes numeric features and assembles them with $q$. We compare depths 3 and 5 with 40 boosting iterations, step size 0.1, a minimum of 25 training examples per child node, subsampling 0.8 and a maximum of 128 feature bins. Every reported configuration selects depth 5. Temperature calibration is applied to the binary probability pair $(1-p,p)$: A selects $T=0.75$ and B selects $T=1$ at each fraction. Unlike the RF survival mixture, this construction does not enforce decreasing probability as $q$ increases.

\subsection{Evaluation protocol}
Adjacent examples share history and future observations, so a random train/test split would mix closely related rainfall episodes. We use chronological partitions (Table~\ref{tab:t2splits}) and remove labels whose target end reaches the following split boundary.
\begin{table}[htbp]\centering\small
\begin{tabular}{lrrp{.29\linewidth}}\toprule
Period & A examples & B examples & Purpose\\\midrule
2017--2021 & 18,688,263 & 27,434,001 & Training population \\
2022 & 6,133,684 & 7,107,898 & Model and cutoff selection \\
2023 & 5,465,591 & 6,719,274 & RF and GBT temperature selection \\
2024 & 5,729,508 & 6,374,476 & Model evaluation \\
\bottomrule\end{tabular}
\caption{Eligible station-time examples after completeness checks and boundary purging. A additionally requires complete published history.}\label{tab:t2splits}
\end{table}

For each classifier in each setting, we compare 15\%, 30\% and 100\% of the eligible training anchors. A deterministic station--time hash (seed 2026) gives nested, label-independent samples. The 30\% samples contain 5,608,769 anchors for A and 8,233,882 for B. Historical frequencies and RF within-bin distributions always use all eligible training anchors; only classifier fitting is subsampled. Selection, calibration and evaluation use the full eligible populations within each setting. Thus the experiment measures classifier fitting-fraction sensitivity; a 15\% fit still uses the full training population for the empirical supports and historical frequencies.

Our primary criterion is mean Brier across the seven fitting thresholds, giving each threshold equal weight. At a particular threshold,
\[
\mathrm{BS}=\frac{1}{N}\sum_{i=1}^{N}(\widehat p_i-z_i)^2,
\qquad z_i=\mathbf1(Y_i>q),
\]
which measures squared probability error; lower is better. We prioritise this score because the task requests probabilities. PR-AUC and ROC-AUC measure event ranking, log loss penalises confident errors, and reliability compares predicted probabilities with observed frequencies. Spark computes AUC with 1,000-bin curve downsampling~\cite{sparkeval}. We examine baseline skill and high-threshold calibration alongside mean Brier to assess performance on rare heavy-rainfall events.

Hyperparameters are selected by 2022 mean Brier. RF/GBT temperature is selected on 2023 from 0.5, 0.75, 1, 1.25, 1.5 and 2; alarm cutoffs use 2022 after applying the selected temperature. RF30 reuses saved reference-configuration weights before calibration and testing. The model/fraction comparison uses 2024 throughout. Later full-training and GBT experiments were developed after the initial 2024 results were inspected. Although 2024 is excluded from fitting and temperature selection, these later comparisons and the final model choice are exploratory rather than an untouched confirmatory test. Execution history and available artifacts are documented in Appendix~\ref{app:experiments}.

\clearpage
\section{Prediction results}
We first assess the value of recent rainfall, then compare model families and training fractions. Calibration and threshold diagnostics explain the remaining prediction errors.

\subsection{Recent rainfall improves forecasts, with smaller gains for heavy rain}

Comparing the 15\% RFs on the same 5,729,508 test examples, A reduces mean Brier from 0.02283 to 0.01613, a \textbf{29.4\% relative reduction} against B. The shared cases cover 89.9\% of B's eligible 2024 set; each forest is fitted on its own eligible training population.
\par
\begin{table}[!ht]\centering\small
\begin{tabular}{rrrrrr}\toprule
$q$ (mm) & Event rate (\%) & Brier A & Brier B & PR-AUC A & PR-AUC B\\\midrule
0 & 6.398 & 0.03798 & 0.05924 & 0.6091 & 0.1097 \\
0.2 & 4.892 & 0.03090 & 0.04614 & 0.5582 & 0.0881 \\
1 & 2.579 & 0.01870 & 0.02499 & 0.4215 & 0.0484 \\
2 & 1.686 & 0.01349 & 0.01652 & 0.3155 & 0.0329 \\
5 & 0.829 & 0.00732 & 0.00820 & 0.1885 & 0.0176 \\
10 & 0.378 & 0.00354 & 0.00376 & 0.1020 & 0.0085 \\
20 & 0.100 & 0.00098 & 0.00100 & 0.0304 & 0.0029 \\
\bottomrule
\end{tabular}
\caption{The 15\% RFs evaluated on identical 2024 examples, with event frequency at each rainfall threshold.}\label{tab:t2common}
\end{table}
\par
The benefit is greatest for any positive rain (35.9\% Brier reduction) and falls to 1.7\% above 20 mm (Table~\ref{tab:t2common}, Figure~\ref{fig:t2common}). Recent history therefore adds useful predictive information in this holdout, but its strongest benefit is at lower rainfall thresholds.
\plot{18_common_comparison}{The 15\% RF comparison on shared 2024 examples. Left: relative Brier reduction with conditional 95\% block-bootstrap intervals. Right: PR-AUC and event prevalence on a logarithmic scale. Rainfall thresholds are displayed as equally spaced categories.}{fig:t2common}
To account for overlapping labels and shared weather across stations, we form 2,000 bootstrap replicates by resampling circular blocks of seven consecutive calendar days, retaining all stations and thresholds together. The 95\% interval for the mean relative Brier reduction is 25.4--32.9\%, conditional on these fitted models and the observed 2024 network.

Equal threshold weights do not equalise the attainable absolute error reduction. The 0 and 0.2~mm queries account for 77.8\% of A's total absolute Brier improvement over B. At 20~mm there are 5,715 positive, overlapping station-time windows: A's Brier is 0.0009791 versus 0.0009975 for always predicting zero probability, only 1.84\% lower; its skill over historical frequency is 2.18\%. Its PR-AUC of 0.0304 still exceeds the event prevalence of about 0.001. Heavy-rain ranking is useful, while the probability-error gain is modest.

\FloatBarrier
\subsection{Model comparison and uncertainty}
Holding the fitting fraction at 100\%, \textbf{RF has the lowest observed mean Brier among the tested classifier pipelines in both settings}. In B, however, the historical-frequency baseline has lower error than all three classifiers (Table~\ref{tab:t2baseline}).
\par
\begin{table}[!ht]\centering\small
\begin{tabular}{llrrrr}\toprule
Mode & Model & Training & Mean Brier & \shortstack{PR-AUC\\at 1 mm} & \shortstack{Log loss\\at 1 mm}\\\midrule
A & Random forest & 100\% & 0.01614 & 0.4229 & 0.0824 \\
A & Logistic regression & 100\% & 0.01944 & 0.3634 & 0.0875 \\
A & GBT & 100\% & 0.01623 & 0.4282 & 0.0824 \\
A & Historical frequency & 100\% & 0.02270 & 0.0467 & 0.1166 \\
B & Random forest & 100\% & 0.02275 & 0.0518 & 0.1167 \\
B & Logistic regression & 100\% & 0.02301 & 0.0428 & 0.1165 \\
B & GBT & 100\% & 0.02288 & 0.0522 & 0.1194 \\
B & Historical frequency & 100\% & 0.02252 & 0.0516 & 0.1146 \\
\bottomrule
\end{tabular}
\caption{Full-training models evaluated on each mode's eligible 2024 population: 5,729,508 examples for A and 6,374,476 for B. Mean Brier averages the seven original thresholds.}\label{tab:t2baseline}
\end{table}
\par
The ranking metrics explain why model choice depends on the objective. In A, GBT has slightly higher PR-AUC at 1 mm (0.4282 versus 0.4229), but higher mean probability squared error. LR has higher ROC-AUC at 1 mm (0.8611 versus 0.8036), while RF has better PR-AUC, Brier and log loss. Ranking events and assigning accurate probabilities therefore describe different aspects of performance.

The paired uncertainty analysis in Table~\ref{tab:pairedmodels} supports a small RF advantage over GBT in A. In B, the mean difference is less clear: the interval includes zero, and GBT has lower Brier at the two lowest thresholds, whereas RF performs better at higher thresholds. Historical frequency has the lowest mean Brier in B; the full RF's mean error is 1.04\% higher, with conditional interval 0.22--1.79\%. This favours the simpler calendar baseline at the evaluated thresholds.

\begin{table}[htbp]\centering\small
\begin{tabular}{llrr}\toprule
Candidate versus reference & Mode & Reduction (\%) & Conditional 95\% interval\\\midrule
RF 100\% versus GBT 100\% & A & +0.560 & [0.222, 0.923] \\
RF 100\% versus GBT 100\% & B & +0.571 & [-0.229, 1.605] \\
RF 100\% versus RF 15\% & A & -0.036 & [-0.092, 0.019] \\
RF 100\% versus RF 15\% & B & -0.047 & [-0.133, 0.021] \\
\bottomrule\end{tabular}
\caption{Relative reduction in mean 2024 Brier, $100(1-\mathrm{BS}_{candidate}/\mathrm{BS}_{reference})$. Point estimates use observed 2024 scores; 95\% intervals use 2,000 paired circular seven-day block resamples. Positive values favour the candidate. Fractions refer to classifier fitting anchors.}\label{tab:pairedmodels}
\end{table}

\FloatBarrier
\subsection{Larger fitting samples bring little improvement at these settings}
The three-fraction comparison examines sensitivity to the number of classifier fitting anchors at the reported settings. \textbf{The recorded mean Brier changes from 15\% to 100\% are below 0.1\% in magnitude for every classifier and setting.}
\begin{table}[!ht]\centering\small
\begin{tabular}{llrrrr}\toprule
Mode & Model & 15\% & 30\% & 100\% & Improvement (\%)\\\midrule
A & Random forest & 0.01613051 & 0.01613151 & 0.01613638 & -0.036 \\
A & Logistic regression & 0.01944150 & 0.01943988 & 0.01944300 & -0.008 \\
A & GBT & 0.01622929 & 0.01623889 & 0.01622724 & +0.013 \\
B & Random forest & 0.02274285 & 0.02274079 & 0.02275362 & -0.047 \\
B & Logistic regression & 0.02301006 & 0.02301067 & 0.02301076 & -0.003 \\
B & GBT & 0.02290321 & 0.02290187 & 0.02288430 & +0.083 \\
\bottomrule\end{tabular}
\caption{Mean 2024 Brier over seven thresholds on each mode's eligible population. Improvement is the relative Brier reduction from 15\% to 100\%; positive means lower error. The 15\% A GBT value comes from a retained aggregate summary; other cells use metric exports.}\label{tab:trainingsampling}
\end{table}
For RF, the 100\% score is 0.036\% higher in A and 0.047\% higher in B than at 15\%. LR also changes little, while GBT improves slightly at full training. A's 30\% GBT score is worse than its 15\% score, so the observed learning curve is not monotonic. The model ordering is unchanged: RF has lower mean Brier than LR and GBT at every fitting fraction in both settings. The paired RF intervals in Table~\ref{tab:pairedmodels} include zero for both settings. These observations apply to one sampling seed, the tested model capacities and the threshold-average objective; they do not establish equivalence at every rainfall threshold or an end-to-end computing-cost reduction. The detailed learning-curve figure is retained in Appendix~\ref{app:experiments}.

The controlled RF comparison on 2022 provides a complementary check. With fitting settings held fixed, increasing the sample from 15\% to 30\% changes mean Brier by only 0.012\% in A's favour and 0.003\% against B; both conditional intervals include zero. A modest increase in tree depth is more useful for A than doubling the sample at the reference settings. Appendix~\ref{sec:rfsensitivity} gives the individual settings and paired intervals.

Even the 15\% samples contain millions of anchors, many with overlapping windows from the same rainfall episode, but only 2,886 (A) and 4,198 (B) totals above 20 mm.

\FloatBarrier
\subsection{Reference-forest calibration and the alarm tradeoff}
For the 15\% RFs at $q=1$ mm, A distinguishes a wider range of risks than B. Mean predicted risk is 2.38\%, against 2.58\% observed, and the ten-bin weighted absolute calibration error is 0.196 percentage points. Most examples occupy low-probability bins; several higher-probability bins underestimate the observed event frequency (Figure~\ref{fig:t2reliability}).
\plot{20_reliability}{Calibration of the 15\% RFs at 1 mm on shared examples. Upper panels compare bin-average predictions with observed event rates; B is enlarged to show its two occupied bins. Lower panels show sample counts in fixed 0.1-wide probability bins across the full probability range.}{fig:t2reliability}

The rainfall threshold defines the event, while the probability cutoff determines whether to issue an alarm. For the 15\% A RF at $q=1$ mm, changing the probability cutoff from 0.5 to the validation-selected 0.2 raises recall from 24.7\% to 46.4\%, while precision falls from 69.3\% to 47.4\%. Alarm cutoffs maximise F1 on 2022 data after applying the temperature selected on 2023. An always-negative rule already achieves about 97.4\% accuracy here, which is why probability quality and the precision--recall tradeoff are more informative than accuracy.

\FloatBarrier
\subsection{Threshold diagnostics: mean-risk bias and unseen queries}
This comparison uses RF, LR and GBT all fitted at 30\%, evaluated on the same eligible population within each setting. It tests two requirements of the probability interface: plausible risk at large rainfall thresholds, and useful predictions between the fitted query values.

\textbf{Large-accumulation risk.} At the matched 30\% fitting fraction, GBT's mean predicted risk above 20 mm is 0.830\% in A and 2.645\% in B, compared with observed frequencies of 0.100\% and 0.103\%. This overprediction helps explain why event ranking and probability error can favour different models. Figure~\ref{fig:t2thresholdrisk} assesses mean-risk bias: agreement of the overall mean with prevalence alone does not establish calibration across individual risk levels. The retained 15\% RF reliability tables put 99.983\% of A's 20~mm cases and all of B's cases in the first 0.1-wide bin. Finer low-risk reliability bins with event counts and block intervals require rescoring individual predictions; they cannot be reconstructed from these coarse aggregates.
\plot{23_threshold_risk}{Observed frequencies and mean predicted probabilities for matched 30\% models at the seven fitting thresholds, on each mode's eligible 2024 examples. The vertical axis is logarithmic and thresholds are equally spaced categories.}{fig:t2thresholdrisk}

\textbf{Queries between fitted thresholds.} The additional queries at 3, 7.5 and 15 mm fall inside RF amount bins and between the values used to fit LR and GBT. RF has lower Brier at all three query values in both settings (Table~\ref{tab:interior}), supporting its within-bin conversion at these evaluated points.
\begin{table}[!ht]\centering\small
\begin{tabular}{lrrrr}\toprule
Mode & $q$ (mm) & RF Brier & LR Brier & GBT Brier\\\midrule
A & 3 & 0.010546 & 0.012359 & 0.010717 \\
A & 7.5 & 0.005051 & 0.006950 & 0.005132 \\
A & 15 & 0.001854 & 0.003118 & 0.001913 \\
B & 3 & 0.012385 & 0.012510 & 0.012856 \\
B & 7.5 & 0.005534 & 0.005663 & 0.006142 \\
B & 15 & 0.001950 & 0.002002 & 0.002620 \\
\bottomrule\end{tabular}
\caption{2024 results at thresholds absent from the fitting query grid. All three algorithms use 30\% of training anchors and the same eligible test population within each mode.}\label{tab:interior}
\end{table}
\FloatBarrier
GBT also exposes a structural issue with an unconstrained threshold predictor: a larger rainfall threshold can produce a larger exceedance probability. Across the ten evaluated thresholds, the proportion of adjacent query pairs with increasing GBT probability is 0.0491\% at 30\% training and 0.0681\% at 100\% for A, compared with 0.5238\% and 0.6164\% for B. RF preserves the required ordering by construction.

Appendix~\ref{app:prediction} retains threshold-specific baseline skill, monthly skill and feature-importance diagnostics for the 15\% models.

\FloatBarrier
\subsection{Worked probability query and model choice}
Table~\ref{tab:workedquery} illustrates the required interface using archived comparison scores for station S06 at 07:10 SGT on 8 November 2024. A retrieves the twelve published history endpoints from 06:05 to 07:00 and predicts accumulation over 07:10--07:40. Changing only $q$ gives non-increasing probabilities. The realised half-hour was dry despite A assigning 53.63\% probability above 1~mm; this illustrates forecast uncertainty and is not a representative performance sample or a newly executed live query.
\begin{table}[htbp]\centering\small
\begin{tabular}{rrr}\toprule
Threshold $q$ (mm) & A RF probability (\%) & B RF probability (\%)\\\midrule
0 & 82.50 & 6.67 \\
1 & 53.63 & 2.66 \\
10 & 13.01 & 0.43 \\
20 & 4.04 & 0.12 \\
\bottomrule\end{tabular}
\caption{Archived 15\% RF scores for one common station-time example. Both columns predict $P(Y>q)$; the observed $Y$ was 0~mm.}\label{tab:workedquery}
\end{table}

The \textbf{15\% depth-8 A RF is the fully evaluated reference for queries with complete published history}. Increasing its fitting fraction gives no clear mean improvement. Appendix~\ref{sec:rfsensitivity} finds lower 2022 Brier for depth 10 (0.01664745 versus 0.01669750 at 15\%), but that validation-only candidate needs calibration and subsequent evaluation before replacing the reference. For calendar-only risk, historical frequency remains the preferred baseline at the evaluated thresholds. Fine-resolution heavy-rain calibration and chronological confirmation of validation-selected candidates are the next priorities.

\FloatBarrier
% !TEX root = ../project_report.tex
\section{Discussion}

Recent rainfall substantially improves the short-term probability estimates, especially at lower thresholds. Calendar features describe recurring seasonal and daily patterns, while recent observations capture local rainfall persistence. Specifically, RF has the lowest observed threshold-average Brier among the classifier pipelines, with a small positive conditional advantage over GBT in A. GBT ranks 1~mm events competitively, and in B its Brier is lower at the two lowest thresholds. Historical frequency is strongest for calendar-only mean Brier. It is also worth noting that limitations like average risk agreement, coarse-bin calibration and ranking performance still exist.

The fitting-fraction comparison shows small mean changes at the chosen capacities and seed. RF still estimates within-bin distributions from the full training population, and LR/GBT expand each fitting anchor into seven binary examples. The results therefore do not quantify end-to-end savings from using less data. The depth-10 validation improvement is more promising than increasing the reference forest's fitting fraction, but it requires its own calibrated evaluation. Later reuse of 2024 limits confirmatory model-selection claims. Rolling chronological comparisons within the development years and evaluation in a subsequent untouched period would assess stability.

Coverage defines the population described. Task 1 retains 85 of 96 year-months; A retains 89.9\% of B's eligible 2024 examples. Its scores therefore apply to times with complete published history and future labels. At 1~mm, the 15\% A RF improves Brier over its same-setting historical-frequency baseline in every 2024 month (9.8--33.9\%) and at all 70 evaluated stations (22.1--41.9\%). These descriptive subgroup comparisons have no subgroup confidence intervals. Broader assessment should cover other thresholds and years. Neither the eight-year observation record nor a single evaluation year establishes long-term climate or model stability.

We checked the station-catalogue concern directly. Although offline B features take coordinates from observations and the query command uses the frozen training catalogue, all saved catalogues agree and every active 2022--2024 station-month matches them within numerical precision, with no within-month coordinate variation. Thus this policy difference does not change coordinates in the evaluated years. Future relocations require consistent catalogue updates. A live versioned archive would additionally establish which record version was available at each issue time.

\section{Conclusion}
The records show seasonal, diurnal and local rainfall variation across Singapore. Coverage checks and common-station comparisons support the descriptive findings, although changing station coverage and source gaps remain limitations. For Task 2, recent-rainfall RF is a useful evaluated reference, with gains concentrated at lower thresholds and limited heavy-rain probability skill. Historical frequency remains strongest for calendar risk at the evaluated thresholds. The close model and fitting-fraction comparisons support scoped conclusions about the retained pipelines; stronger validation and fine-resolution tail diagnostics are the next steps.

\par\noindent\begin{minipage}{\linewidth}
\section*{AI usage declaration}
OpenAI Codex assisted with code development, cloud execution, and report drafting.
\end{minipage}

\begingroup\small\raggedright
\begin{thebibliography}{9}
\bibitem{data} National Environment Agency. \emph{Historical Rainfall across Singapore}, annual datasets 2017--2024. Data.gov.sg. \url{https://data.gov.sg/collections/2279/view}. Accessed 9 October 2026.
\bibitem{climate} Meteorological Service Singapore. \emph{Climate of Singapore}. \url{https://www.weather.gov.sg/climate-climate-of-singapore/}. Accessed 9 October 2026.
\bibitem{geo} geoBoundaries / Urban Redevelopment Authority. Singapore ADM0 (SGP-ADM0-21272760) and ADM2 (SGP-ADM2-91842731), commit \texttt{9469f09}, built 12 December 2023; based on URA Master Plan 2014. Retained metadata identifies ODbL 1.0 for ADM0 and Singapore Open Data Licence 1.0 for ADM2. \url{https://github.com/wmgeolab/geoBoundaries/tree/9469f09/releaseData/gbOpen/SGP}. Retrieved 8 October 2026.
\bibitem{sparkml} Apache Spark. \emph{Classification and regression}, Spark 3.5.3 documentation. \url{https://archive.apache.org/dist/spark/docs/3.5.3/ml-classification-regression.html}. Accessed 17 October 2026.
\bibitem{sparkeval} Apache Spark. \emph{BinaryClassificationMetrics}, Spark 3.5.3 Java API. \url{https://archive.apache.org/dist/spark/docs/3.5.3/api/java/org/apache/spark/mllib/evaluation/BinaryClassificationMetrics.html}. Accessed 17 October 2026.
\end{thebibliography}
\endgroup
% !TEX root = ../project_report.tex
\clearpage
\appendix
\numberwithin{figure}{section}
\numberwithin{table}{section}
\setcounter{figure}{0}
\setcounter{table}{0}
\renewcommand{\thefigure}{\Alph{section}.\arabic{figure}}
\section{Distribution and source-gap evidence}\label{app:data}
These supplementary figures separate the distribution of measured rain from the absence of measurements. They support the data-quality discussion without treating an unrecorded interval as dry weather.
\plot{02_positive_distribution}{Distribution of positive five-minute rainfall across 2017--2024. Bars give the share of positive observations in log-spaced bins; the horizontal axis is logarithmic. Recorded zeros are excluded. The long upper tail motivates evaluating rare accumulations separately.}{fig:appdistribution}
\plot{04_source_gap}{Daily network-wide raw and cleaned record counts in August 2018. The absence of observations on 4--31 August is already present in the source file. The cleaning process therefore does not explain this gap, and this figure does not identify its upstream cause.}{fig:appgap}
\FloatBarrier

\clearpage
\section{Monthly trajectories and annual support}\label{app:temporal}
The monthly trajectories retain usable periods individually. The annual plot applies the stricter full-year rule, showing why monthly evidence is more extensive than annual evidence.
\plot{07_monthly_trends}{Mean observed monthly station totals, with a separate trajectory for each calendar month. Missing year-months break the lines. Vertical ranges differ across panels; use the main monthly heatmap to compare absolute amounts between months. Station composition can change between points, so these are descriptive network trajectories.}{fig:apptrajectories}
\plot{08_annual_totals}{Observed calendar-year totals at eligible stations. Points show station totals, coral diamonds their mean, and $n$ the contributing station count. Grey columns have no station meeting annual coverage of at least 95\% and coverage of at least 90\% in every month. The three supported years use different station sets and do not establish a national wettest-year ranking.}{fig:appannual}
\FloatBarrier

\clearpage
\section{Daily rainfall and coverage sensitivity}\label{app:robustness}
Daily upper-tail amounts and wet-day frequencies describe different aspects of rainfall. The sensitivity curves show how changing the completeness rule affects seasonal summaries, while retaining the limitations of the observed network.
\plot{11_daily_tails}{Monthly daily-rainfall P99 (upper) and wet station-day percentage (lower). Included days have at least 95\% daily coverage within eligible station-months. Quantiles include dry days; wet means observed daily rainfall above zero. P99 is an approximate Spark quantile of station-days, not an independent-storm return level. Grey cells lack qualifying support.}{fig:apptails}
\plot{12_sensitivity}{Seasonal mean amount and wet-interval frequency under 80\%, 90\% and 95\% monthly coverage cutoffs. Each curve uses its own eligible stations and equally weighted available years. Similar curves support the stability of these summaries under the tested cutoffs; they do not establish random missingness.}{fig:appsensitivity}
\FloatBarrier

\clearpage
\section{Additional spatial views}\label{app:spatial}
The twelve-month view complements the main December comparison. Regional colours summarise contributing gauges and should be read together with their locations and coverage.
\plot{14_spatial_year_2024}{All twelve months of 2024. Colours show observed monthly totals at qualified stations, and $n$ gives their number. The original gallery's global monthly colour scale is retained; it differs from the December-only scale in the main text. The eligible network changes by month, and no rainfall is interpolated between gauges.}{fig:appspatialseason}
\plot{15_regions_month_12}{December planning-area station means across eight years. Each coloured polygon shows the arithmetic mean of its assigned eligible stations; dots show contributing locations. Grey means no supporting station, not zero rainfall. The original regional scale is retained. These are gauge summaries, not area-weighted rainfall estimates; station counts and placements affect their representativeness.}{fig:appregions}
\FloatBarrier

\clearpage
\section{Additional predictive diagnostics}\label{app:prediction}
These figures describe the 15\% RF and logistic models on the 2024 test set. They show how performance varies by rainfall threshold and month, alongside the RF's feature use.
\plot{19_baseline_skill}{Brier skill relative to each mode's historical-frequency baseline on its own eligible 2024 test population. Positive values indicate lower error than the baseline; zero indicates equal error. A's random forest improves on the baseline at every evaluated threshold, whereas A's logistic regression deteriorates at larger thresholds. B's forest has small, threshold-dependent gains and losses. The two panels use different eligible populations.}{fig:appbaselines}
\plot{21_diagnostics}{Left: A's monthly Brier skill relative to historical frequency at 1 mm in 2024. Right: its six largest random-forest impurity importances. The 15-minute total has the largest impurity importance. These importances describe split use in the fitted forest; correlated rainfall summaries share information and the bars do not measure independent causal effects.}{fig:appdiagnostics}
\FloatBarrier

\clearpage
\section{Model scores and experiment coverage}\label{app:experiments}
Table~\ref{tab:allmodels} brings the recorded model and training-fraction results together. Mean Brier and mean log loss give equal weight to the seven primary thresholds; PR-AUC is shown at 1 mm. Historical frequency uses all eligible training examples in every run. The fitting fraction refers to classifier anchors before LR/GBT expansion into seven threshold-specific examples; it does not subsample the test set.
\begin{table}[htbp]\centering\small
\begin{tabular}{llrrrr}\toprule
Mode & Model & Fitting sample & Mean Brier & Mean log loss & PR-AUC at 1 mm\\\midrule
A & Random forest & 15\% & 0.01613051 & 0.070004 & 0.4218 \\
A & Random forest & 30\% & 0.01613151 & 0.070021 & 0.4225 \\
A & Random forest & 100\% & 0.01613638 & 0.070016 & 0.4229 \\
A & Logistic regression & 15\% & 0.01944150 & 0.082787 & 0.3629 \\
A & Logistic regression & 30\% & 0.01943988 & 0.082783 & 0.3632 \\
A & Logistic regression & 100\% & 0.01944300 & 0.082788 & 0.3634 \\
A & GBT & 15\% & 0.01622929 & 0.072167 & --- \\
A & GBT & 30\% & 0.01623889 & 0.072197 & 0.4203 \\
A & GBT & 100\% & 0.01622724 & 0.072151 & 0.4282 \\
A & Historical frequency & 100\% & 0.02269628 & 0.099923 & 0.0467 \\
B & Random forest & 15\% & 0.02274285 & 0.100015 & 0.0494 \\
B & Random forest & 30\% & 0.02274079 & 0.099983 & 0.0488 \\
B & Random forest & 100\% & 0.02275362 & 0.100170 & 0.0518 \\
B & Logistic regression & 15\% & 0.02301006 & 0.104528 & 0.0428 \\
B & Logistic regression & 30\% & 0.02301067 & 0.104539 & 0.0428 \\
B & Logistic regression & 100\% & 0.02301076 & 0.104540 & 0.0428 \\
B & GBT & 15\% & 0.02290321 & 0.107891 & 0.0511 \\
B & GBT & 30\% & 0.02290187 & 0.107892 & 0.0521 \\
B & GBT & 100\% & 0.02288430 & 0.107829 & 0.0522 \\
B & Historical frequency & 100\% & 0.02252051 & 0.098657 & 0.0516 \\
\bottomrule\end{tabular}
\caption{2024 results on each mode's eligible population: 5,729,508 anchors for A and 6,374,476 for B. The 15\% A GBT row uses its recorded aggregate summary; the 15\% B GBT row uses per-threshold exports. All other rows use saved metric CSVs.}\label{tab:allmodels}
\end{table}

\plot{22_training_fraction}{Relative improvement in mean 2024 Brier over each model's 15\% result; positive means lower error. The open A GBT marker identifies its aggregate-summary reference. Small point differences should be read with the paired intervals in the main results.}{fig:t2sampling}

The table covers all three fitting fractions for RF, LR and GBT. RF 30\% reuses the standard depth-8, 40-tree forests; LR 30\% uses separate fits. All classifier rows retain pipelines and metric exports except GBT 15\%: A has an aggregate summary and B has per-threshold text exports. RF/LR columns inside GBT exports rescore the 15\% reference models. The Part 2 repository guide records the source paths.

\begin{table}[htbp]\centering\small
\begin{tabular}{>{\raggedright\arraybackslash}p{.27\linewidth}>{\raggedright\arraybackslash}p{.32\linewidth}>{\raggedright\arraybackslash}p{.31\linewidth}}\toprule
Execution batch & Fitted or reused models & Recorded evaluation\\\midrule
Initial experiment & RF and LR at 15\%, A and B & 2022 selection; RF calibration on 2023; 2024 test\\
RF sensitivity & RF at 15\% and 30\%; one-factor tuning & 2022 validation\\
Full training & RF and LR at 100\%, A and B & 2022 selection; RF calibration on 2023; 2024 test\\
GBT experiment & GBT at 15\%, 30\%, 100\%, A and B & 2022 selection; 2023 calibration; 2024 test\\
LR 30\% completion & Newly fitted LR at 30\%, A and B & 2022 selection; 2024 test\\
RF 30\% completion & Reused standard RF 30\% weights, A and B & 2023 calibration; 2022 cutoff selection; 2024 test\\\bottomrule
\end{tabular}
\caption{Execution provenance for the unified comparison. All classifier fitting uses 2017--2021. Later batches reuse 2024 after the initial results were available; they are not a new independent test year.}\label{tab:experiment_batches}
\end{table}
\FloatBarrier

\clearpage
% !TEX root = ../project_report.tex
\section{Random-forest tuning on 2022 validation}\label{sec:rfsensitivity}

This experiment varies the training fraction and RF settings on the 2022 validation population. The reference forest uses 40 trees, depth 8, minimum leaf size 25, square-root feature selection and subsampling of 0.8. A deterministic station--time hash with seed 2026 gives nested 15\% and 30\% samples. The 30\% samples contain 5,608,769 anchors for A and 8,233,882 for B. Every candidate within a mode uses the same validation anchors, training-derived within-bin distributions and temperature $T=1$.

At 15\%, the one-factor comparisons increase depth to 10, increase the number of trees to 80, increase minimum leaf size to 50, or use half the features at each split. The lowest-Brier configuration at 15\% is then fitted at 30\%. Table~\ref{tab:rfsensitivity} gives all candidate scores. These are one-factor comparisons, rather than a search over all parameter combinations.

\par
\begin{table}[htbp]\centering\small
\begin{tabular}{lrrr}\toprule
Configuration & Fitting sample & A mean Brier & B mean Brier\\\midrule
Reference & 15\% & 0.01669750 & 0.02332733 \\
Reference & 30\% & 0.01669544 & 0.02332800 \\
Depth 10 & 15\% & 0.01664745 & 0.02331832 \\
80 trees & 15\% & 0.01669631 & 0.02332818 \\
Minimum leaf size 50 & 15\% & 0.01671066 & 0.02333420 \\
Half-feature selection & 15\% & 0.01668206 & 0.02331674 \\
Selected setting & 30\% & 0.01663765 & 0.02330797 \\
\bottomrule\end{tabular}
\caption{Mean Brier on the full 2022 validation set, averaged over seven rainfall thresholds. Lower is better. Changes at 15\% vary one setting from the reference. The selected 30\% row uses depth 10 for A and half-feature selection for B; it is not the same configuration for both modes.}
\label{tab:rfsensitivity}\end{table}\par


At the reference settings, increasing the sample to 30\% reduces A's mean Brier from 0.01669750 to 0.01669544, a relative improvement of 0.012\%. B's mean Brier increases from 0.02332733 to 0.02332800, a relative deterioration of 0.003\%. The corresponding 95\% intervals for improvement are $[-0.014\%,\,0.039\%]$ for A and $[-0.070\%,\,0.047\%]$ for B.

Depth has a larger effect for A. Moving from depth 8 to 10 at 15\% improves mean Brier by 0.300\%, with interval $[0.213\%,\,0.389\%]$. Fitting that depth-10 configuration at 30\% adds a further 0.059\% improvement, with interval $[0.036\%,\,0.082\%]$, and reaches 0.01663765. B's best tested 15\% setting uses half the features; its 30\% version reaches 0.02330797. Both B intervals for the configuration and subsequent sample changes cross zero. At the reference settings, A's Brier improves by 0.103\% at 10 mm and 0.072\% at 20 mm.

Intervals use 2,000 paired circular resamples of seven-day blocks, retaining all stations and thresholds together. They describe validation-day variation for the fitted models and seed used here. Configuration selection uses the same validation year. These Brier comparisons complement the 2024 model and sample-size results in the main text.

The reference 30\% forests in this table were subsequently reused for the 2024 comparison in Appendix~\ref{app:experiments}, with temperature selected on 2023 ($T=1$ in both settings). The selected-setting 30\% forests remain validation-only candidates. The scores here describe 2022, while the main sample-size table describes 2024, so the numerical gains refer to different evaluation years.

\FloatBarrier

\end{document}
